\chapter*{Colophon}
\addcontentsline{toc}{chapter}{Colophon}
\textbf{Authorship and copyright.} Blendity is copyright Emiliano Manalo and contributors. Its source code, tests,
documentation and these two books were written with Anthropic's Claude through Claude Code \cite{ClaudeCode2026}, under
a human's direction and review: the user decides what Blendity is for, approves each task's specification and decision
record before code is written, reviews every pull request and merges it. The model's outputs belong to the user who
directed them under Anthropic's terms \cite{AnthropicTerms2026}, and are licensed to the extent copyright subsists in
them (the human selection, direction, review and editing); the \code{Co-Authored-By} trailer on every commit is
attribution, not a licence grant. No third-party code is claimed as the project's own: what is copied is listed in
\code{licenses/README.md} with its licence, what is followed is credited in \code{CREDITS.md} and the References
chapter, and the whole story is in \code{docs/AI\_DISCLOSURE.md}.

\textbf{Licence.} The source is GPL-2.0-or-later \cite{GPLv2}, the same as Blender's \cite{Blender2026}; the
executables, which combine it with Apache-2.0 libraries \cite{Apache2}, are distributed under GPL-3.0-or-later
\cite{GPLv3}. The scenes, meshes, textures, renders and animations made with Blendity are their maker's: the GPL covers
the program, not its output.

\textbf{How the books are made.} The sources are LaTeX under \code{docs/book/}, typeset with LuaLaTeX in the DejaVu
faces (the editor's own) and biblatex with biber for the reference list. \code{tools/make\_books.py} generates the
reference list and the index tables from the project's own files, takes the figures headless from a scratch project,
builds both PDFs and checks that nothing has drifted. The release workflow builds them again on every tag and attaches
them to the release.

\textbf{Accessibility and readability.} The books follow the guidelines the editor itself is measured against
(\adr{0085}), applied to documents: Nielsen's tenth heuristic, help and documentation \cite{Nielsen1994Enhancing,
Nielsen1993} (easy to search: a table of contents, bookmarks, a list of figures, a reference list with back-references
and three indexes; focused on the user's task: the user guide is the work stage by stage with numbered steps and worked
examples, the architecture book the questions a maintainer asks; not too large: each chapter ends in a pointer to the
Learn lesson that goes deeper); Carroll's minimalist instruction \cite{Carroll1990} (start on a real task, keep the
reading short, support error recovery: a troubleshooting chapter and the ``when something refuses'' list); and WCAG 2.1
\cite{W3C2018Web} as PDF/UA applies it \cite{PDFUA2024}: every text colour, link and box title reads at 4.5:1 or more
(the ratios are in \code{blendity-common.sty}), links are told apart by form as well as colour (bracketed citation
numbers, ``chapter 3'', monospace URLs), every figure carries alt text and a caption that says what it shows, the
diagrams' meaning is in their captions and the text around them, the document's language and metadata are declared
and its headings are the PDF's bookmarks (structure tags for a screen reader are the one open item: LaTeX's tagging
project does not yet support the packages the books use, so the captions carry what the alt text would), the body
type is about 10 points of a face made for screens with 80 characters a line, and nothing is conveyed by colour
alone. Keys, menus and controls are
typeset consistently: \key{Ctrl+K} is a key, \menu{Window > Layouts} a menu path, \ui{Generate Lighting} a control's
label, \code{make\_books.py} a file or command.

\textbf{Version.} Blendity 0.13.0, November 2026, from the \code{main} branch of
\url{https://github.com/EmilianoManaloIV/blendity}.
